\documentclass[11pt,a4paper]{article}

\usepackage[a4paper,margin=25mm]{geometry}
\usepackage{fontspec}
\usepackage{polyglossia}
\setmainlanguage{vietnamese}
\IfFontExistsTF{Times New Roman}{\setmainfont{Times New Roman}}{\setmainfont{Latin Modern Roman}}
\IfFontExistsTF{Arial}{\setsansfont{Arial}}{\setsansfont{Latin Modern Sans}}
\IfFontExistsTF{Courier New}{\setmonofont{Courier New}}{\setmonofont{Latin Modern Mono}}
\usepackage{amsmath,amssymb}
\usepackage{array}
\usepackage{booktabs}
\usepackage{enumitem}
\usepackage{hyperref}
\usepackage{listings}
\usepackage{xcolor}

\hypersetup{
  colorlinks=true,
  linkcolor=blue!60!black,
  urlcolor=blue!60!black,
  pdftitle={Bao cao luc va moment cua Main Rotor va Tail Rotor},
  pdfauthor={Codex}
}

\lstset{
  basicstyle=\ttfamily\small,
  frame=single,
  breaklines=true,
  columns=fullflexible
}

\newcommand{\code}[1]{\texttt{\detokenize{#1}}}
\newcommand{\vect}[1]{\boldsymbol{#1}}
\newcommand{\R}{\mathbb{R}}

\title{Báo Cáo Cách Tính Lực và Moment\\từ mã C của Main Rotor và Tail Rotor}
\author{Phân tích trực tiếp từ \code{MainRotor_ert_rtw} và \code{TailRotor_ert_rtw}}
\date{06/04/2026}

\begin{document}
\maketitle
\tableofcontents

\section{Phạm vi}

Báo cáo này chỉ dựa trên mã C sinh tự động trong bốn tệp:
\begin{itemize}[leftmargin=1.8em]
  \item \code{MainRotor_ert_rtw/MainRotor.c}
  \item \code{MainRotor_ert_rtw/MainRotor_data.c}
  \item \code{TailRotor_ert_rtw/TailRotor.c}
  \item \code{TailRotor_ert_rtw/TailRotor_data.c}
\end{itemize}

Mục tiêu là mô tả đúng cách lực và moment được tính ở đầu ra:
\[
  \vect{F}_m,\ \vect{M}_m \quad \text{cho main rotor}, \qquad
  \vect{F}_t,\ \vect{M}_t \quad \text{cho tail rotor}.
\]

Tôi không dùng mô hình ground effect cũ và không dựa vào hình Simulink; mọi công
thức dưới đây đều được rút trực tiếp từ mã C đã sinh.

\section{Giao diện hai hàm tính toán}

\subsection{Main rotor}

\begin{lstlisting}[language=C]
void MainRotor_step(
  RT_MODEL_MainRotor_T *const MainRotor_M,
  real_T SpeedM,
  real_T CollectiveM,
  real_T CyclicLat,
  real_T CyclicLon,
  real_T rho,
  real_T wb[3],
  real_T Vb[3],
  real_T Fm[3],
  real_T Mm[3]);
\end{lstlisting}

\begin{center}
\begin{tabular}{@{}ll@{}}
\toprule
Biến & Ý nghĩa \\
\midrule
\code{SpeedM} & tốc độ góc rotor chính, ký hiệu \(\Omega_m\) \\
\code{CollectiveM} & collective pitch của main rotor \\
\code{CyclicLat}, \code{CyclicLon} & cyclic pitch ngang và dọc \\
\code{rho} & mật độ không khí \(\rho\) \\
\code{wb[3]} & tốc độ góc thân \\
\code{Vb[3]} & vận tốc thân \\
\code{Fm[3]} & lực đầu ra của main rotor \\
\code{Mm[3]} & moment đầu ra của main rotor \\
\bottomrule
\end{tabular}
\end{center}

\subsection{Tail rotor}

\begin{lstlisting}[language=C]
void TailRotor_step(
  RT_MODEL_TailRotor_T *const TailRotor_M,
  real_T SpeedT,
  real_T CollectiveT,
  real_T rho,
  real_T Ft[3],
  real_T Mt[3]);
\end{lstlisting}

\begin{center}
\begin{tabular}{@{}ll@{}}
\toprule
Biến & Ý nghĩa \\
\midrule
\code{SpeedT} & tốc độ góc tail rotor, ký hiệu \(\Omega_t\) \\
\code{CollectiveT} & collective pitch của tail rotor \\
\code{rho} & mật độ không khí \(\rho\) \\
\code{Ft[3]} & lực đầu ra của tail rotor \\
\code{Mt[3]} & moment đầu ra của tail rotor \\
\bottomrule
\end{tabular}
\end{center}

\section{Phần chung: cách mã C tính CT và CQ}

Hai rotor dùng cùng một lõi blade-element/momentum được rút ra từ code.
Khác nhau nằm ở tham số hình học và ở phần biến đổi phương lực sau khi đã có độ lớn lực.

\subsection{Pitch hiệu dụng và solidity}

Với mỗi rotor, mã C trước tiên tính:
\begin{equation}
  \bar{\theta}
  = \frac{2}{3}\left(\theta_{\mathrm{root}} + \theta_{\mathrm{cmd}}\right)
    + \frac{1}{2}\theta_{\mathrm{twist}},
\end{equation}
và
\begin{equation}
  \sigma = \frac{N_b c}{\pi R}.
\end{equation}

\subsection{Rời rạc bán kính cánh}

Mỗi rotor chia span thành 100 đoạn đều từ
\begin{equation}
  r_0 = \frac{e}{R}
\end{equation}
đến \(1\), với
\begin{equation}
  \Delta r = \frac{1-e/R}{100}.
\end{equation}
Tại mỗi midpoint \(r_i\), code lặp Newton để tìm một inflow cục bộ \(\lambda_i\).

\subsection{Phương trình Newton trong code}

Từ biểu thức cập nhật trong \code{MainRotor.c} và \code{TailRotor.c}, phương trình
mà code đang giải tại mỗi station có thể viết thành
\begin{equation}
  4 F_i(\lambda_i)\lambda_i^2
  - \frac{\sigma a}{2}\left(\bar{\theta}-\lambda_i\right) = 0,
  \label{eq:newton-balance}
\end{equation}
Phương trình này được dùng để tính inflow vì \(\lambda_i\) chính là ẩn duy nhất
chưa biết tại mỗi radial station. Vế thứ nhất
\(4F_i(\lambda_i)\lambda_i^2\) là phần thrust theo momentum theory có kể đến
tip-loss, còn vế thứ hai
\(\frac{\sigma a}{2}\left(\bar{\theta}-\lambda_i\right)\) là phần thrust theo
blade-element theory với giả thiết góc tấn nhỏ. Khi hai vế này bằng nhau, code
đã tìm được giá trị inflow cục bộ sao cho tải khí động của phần tử cánh phù hợp
với induced flow mà rotor tự tạo ra. Nói ngắn gọn, phương trình (5) không phải
là công thức ``đặt sẵn'' cho inflow, mà là phương trình cân bằng cần phải giải
ra để lấy inflow.

Nếu dùng công thức inflow đơn giản kiểu hover đều, ví dụ
\(\lambda \approx \sqrt{C_T/2}\), thì sẽ chỉ thu được một giá trị trung bình cho
toàn đĩa rotor. Trong khi đó code đang cần \(\lambda_i\) theo từng station để
tích phân \(C_T\) và \(C_Q\), đồng thời vẫn giữ được ảnh hưởng của tip-loss và
pitch hiệu dụng trong mô hình sinh tự động.

trong đó hệ số tip-loss là
\begin{equation}
  F_i(\lambda_i)
  = \frac{2}{\pi}
    \cos^{-1}\!\left[
      \exp\!\left(
        -\frac{N_b}{2}\frac{1-r_i}{\lambda_i}
      \right)
    \right].
  \label{eq:tip-loss}
\end{equation}

\paragraph{Lưu ý quan trọng về mã sinh tự động.}
Điều kiện dừng trong cả hai file đều có dạng
\code{fabs(x - x) >= 1.0E-6}, nên luôn bằng 0. Vì vậy vòng Newton thực tế
không dừng theo hội tụ mà chạy cố định đúng 100 bước mỗi station.

\subsection{Tích phân hệ số thrust và torque}

Sau khi có \(\lambda_i\), code cộng dồn:
\begin{align}
  C_T &\approx \sum_{i=1}^{100} 4F_i\lambda_i^2 r_i \Delta r, \\
  C_{Q_i} &\approx \sum_{i=1}^{100} 4F_i\lambda_i^3 r_i \Delta r.
\end{align}

Sau đó hệ số torque tổng là
\begin{equation}
  C_Q = C_{Q_i} + \frac{\sigma c_{d0}}{8}.
\end{equation}

Trong code:
\begin{itemize}[leftmargin=1.8em]
  \item \code{Rotor1_o3}, \code{Rotor6_o3} là \(C_T\)
  \item \code{Rotor1_o4}, \code{Rotor6_o4} là \(C_Q\)
\end{itemize}

\subsection{Độ lớn thrust và torque vật lý}

Với mỗi rotor:
\begin{align}
  T &= C_T\,\rho\,\pi R^2(\Omega R)^2, \\
  Q &= C_Q\,\rho\,\pi R^5\,\Omega |\Omega|.
\end{align}

Trong code, \(T\) dùng bình phương \((\Omega R)^2\) nên luôn là độ lớn dương,
còn \(Q\) dùng \(\Omega|\Omega|\) nên giữ dấu quay của rotor.

\section{Main Rotor}

\subsection{Tham số của main rotor}

\begin{center}
\begin{tabular}{@{}lll@{}}
\toprule
Ký hiệu & Giá trị & Ý nghĩa \\
\midrule
\(R_m\) & \(7.35\) & bán kính rotor chính \\
\(N_{b,m}\) & \(2\) & số cánh \\
\(c_m\) & \(0.53\) & chord \\
\(e_m\) & \(0\) & hinge offset \\
\(a_m\) & \(5.5\) & lift-curve slope \\
\(c_{d0,m}\) & \(0.009\) & profile drag coefficient \\
\(\gamma_m\) & \(10\) & Lock number / tham số flap trong block \\
\(\theta_{\mathrm{root},m}\) & \(0.1493\) & root pitch \\
\(\theta_{\mathrm{twist},m}\) & \(-0.1745\) & blade twist \\
\(\vect{r}_m\) & \([ -0.165,\ 0,\ -2.07]^T\) & vector từ mốc thân tới hub main rotor \\
\bottomrule
\end{tabular}
\end{center}

\subsection{Độ lớn thrust của main rotor}

Sau khi hoàn tất phần tích phân BEMT, code tạo độ lớn thrust:
\begin{equation}
  T_m = C_{T,m}\,\rho\,\pi R_m^2(\Omega_m R_m)^2.
\end{equation}

\subsection{Hướng của lực main rotor}

Khác với tail rotor, main rotor không dùng hướng cố định. Ở đoạn
\code{MainRotor.c} từ khoảng dòng 198 đến 288, code tính hai góc trung gian
\(\alpha_m\) và \(\beta_m\) rồi lưu vào \code{tmp[0]} và \code{tmp[1]}.

Hai góc này phụ thuộc vào:
\begin{itemize}[leftmargin=1.8em]
  \item vận tốc thân \(\vect{V}_b\)
  \item tốc độ góc thân \(\vect{\omega}_b\) từ \code{wb[3]}
  \item hai lệnh cyclic \code{CyclicLat}, \code{CyclicLon}
  \item tham số \(\gamma_m\), \(\sigma_m\), \(C_{T,m}\), \(R_m\)
\end{itemize}

Nói ngắn gọn, code dùng chúng để xác định độ nghiêng của đĩa rotor trước khi
chiếu thrust vào hệ trục thân.

Sau đó lực đầu ra được ghi đúng theo:
\begin{equation}
  \vect{F}_m
  =
  \begin{bmatrix}
    -T_m\cos\beta_m\sin\alpha_m \\
    \phantom{-}T_m\sin\beta_m \\
    -T_m\cos\beta_m\cos\alpha_m
  \end{bmatrix}.
  \label{eq:main-force}
\end{equation}

Tức là:
\begin{align}
  F_{m,x} &= -T_m\cos\beta_m\sin\alpha_m, \\
  F_{m,y} &= \phantom{-}T_m\sin\beta_m, \\
  F_{m,z} &= -T_m\cos\beta_m\cos\alpha_m.
\end{align}

Nếu bỏ ảnh hưởng của \(\vect{V}_b\), \(\vect{\omega}_b\) và cyclic pitch,
thì \(\alpha_m=\beta_m=0\), khi đó:
\[
  \vect{F}_m = [0,\ 0,\ -T_m]^T.
\]

\subsection{Moment của main rotor}

Code tính moment theo hai phần:
\begin{enumerate}[leftmargin=1.8em]
  \item moment do lực đặt lệch tâm: \(\vect{r}_m \times \vect{F}_m\)
  \item moment cản trục rotor: \([0,\ 0,\ -Q_m]^T\)
\end{enumerate}

Vì vậy:
\begin{equation}
  \vect{M}_m
  =
  \vect{r}_m \times \vect{F}_m
  +
  \begin{bmatrix}
    0 \\ 0 \\ -Q_m
  \end{bmatrix}.
  \label{eq:main-moment}
\end{equation}

Khai triển với \(\vect{r}_m=[x_m,\ y_m,\ z_m]^T=[-0.165,\ 0,\ -2.07]^T\):
\begin{align}
  M_{m,x} &= y_m F_{m,z} - z_m F_{m,y} = 2.07\,F_{m,y}, \\
  M_{m,y} &= z_m F_{m,x} - x_m F_{m,z}
          = -2.07\,F_{m,x} + 0.165\,F_{m,z}, \\
  M_{m,z} &= x_m F_{m,y} - y_m F_{m,x} - Q_m
          = -0.165\,F_{m,y} - Q_m.
\end{align}

\section{Tail Rotor}

\subsection{Tham số của tail rotor}

\begin{center}
\begin{tabular}{@{}lll@{}}
\toprule
Ký hiệu & Giá trị & Ý nghĩa \\
\midrule
\(R_t\) & \(1.29\) & bán kính tail rotor \\
\(N_{b,t}\) & \(2\) & số cánh \\
\(c_t\) & \(0.2133\) & chord \\
\(e_t\) & \(0\) & hinge offset \\
\(a_t\) & \(5.5\) & lift-curve slope \\
\(c_{d0,t}\) & \(0.002\) & profile drag coefficient \\
\(\theta_{\mathrm{root},t}\) & \(0.015\) & root pitch \\
\(\theta_{\mathrm{twist},t}\) & \(0\) & blade twist \\
\(\vect{r}_t\) & \([ -8.53,\ 0,\ -2.10]^T\) & vector từ mốc thân tới hub tail rotor \\
\bottomrule
\end{tabular}
\end{center}

\subsection{Độ lớn thrust và torque của tail rotor}

Phần BEMT của tail rotor giống main rotor, chỉ khác tham số số học.
Do đó:
\begin{align}
  T_t &= C_{T,t}\,\rho\,\pi R_t^2(\Omega_t R_t)^2, \\
  Q_t &= C_{Q,t}\,\rho\,\pi R_t^5\,\Omega_t |\Omega_t|.
\end{align}

\subsection{Lực tail rotor trong hệ cục bộ rotor}

Ngay sau block rotor, code tạo force và torque cục bộ:
\begin{equation}
  \vect{F}_{t,\mathrm{local}}
  =
  \begin{bmatrix}
    0 \\ 0 \\ -T_t
  \end{bmatrix},
  \qquad
  \vect{M}_{t,\mathrm{local}}
  =
  \begin{bmatrix}
    0 \\ 0 \\ -Q_t
  \end{bmatrix}.
\end{equation}

\subsection{Phép quay cố định của tail rotor}

Không giống main rotor, tail rotor không có \code{Vb}, \code{wb}, hay cyclic pitch.
Hướng lực của nó được cố định bằng một phép quay \(90^\circ\) quanh trục \(x\).

Code làm việc đó bằng:
\begin{itemize}[leftmargin=1.8em]
  \item \code{Constant3 = 90 deg}
  \item đổi sang rad ở dòng unit conversion
  \item tạo DCM với bộ góc \((\phi,\theta,\psi)=(90^\circ,0,0)\)
\end{itemize}

Ma trận quay tương đương là:
\begin{equation}
  R_x\!\left(\frac{\pi}{2}\right)
  =
  \begin{bmatrix}
    1 & 0 & 0 \\
    0 & 0 & -1 \\
    0 & 1 & 0
  \end{bmatrix}.
\end{equation}

Vì vậy:
\begin{equation}
  \vect{F}_t
  = R_x\!\left(\frac{\pi}{2}\right)\vect{F}_{t,\mathrm{local}}
  =
  \begin{bmatrix}
    0 \\ T_t \\ 0
  \end{bmatrix},
\end{equation}
và
\begin{equation}
  \vect{M}_{Q,t}
  = R_x\!\left(\frac{\pi}{2}\right)\vect{M}_{t,\mathrm{local}}
  =
  \begin{bmatrix}
    0 \\ Q_t \\ 0
  \end{bmatrix}.
\end{equation}

\subsection{Moment của tail rotor}

Tail rotor cũng cộng hai phần:
\begin{equation}
  \vect{M}_t
  =
  \vect{r}_t \times \vect{F}_t + \vect{M}_{Q,t}.
\end{equation}

Với \(\vect{r}_t=[x_t,\ y_t,\ z_t]^T=[-8.53,\ 0,\ -2.10]^T\)
và \(\vect{F}_t=[0,\ T_t,\ 0]^T\), ta được:
\begin{align}
  M_{t,x} &= y_t F_{t,z} - z_t F_{t,y}
          = 2.10\,T_t, \\
  M_{t,y} &= z_t F_{t,x} - x_t F_{t,z} + Q_t
          = Q_t, \\
  M_{t,z} &= x_t F_{t,y} - y_t F_{t,x}
          = -8.53\,T_t.
\end{align}

Do đó đầu ra tail rotor có dạng rất gọn:
\begin{equation}
  \vect{F}_t =
  \begin{bmatrix}
    0 \\ T_t \\ 0
  \end{bmatrix},
  \qquad
  \vect{M}_t =
  \begin{bmatrix}
    2.10\,T_t \\
    Q_t \\
    -8.53\,T_t
  \end{bmatrix}.
  \label{eq:tail-final}
\end{equation}

\section{So sánh ngắn giữa hai rotor}

\begin{center}
\begin{tabular}{@{}p{3.2cm}p{5.2cm}p{5.2cm}@{}}
\toprule
Đặc điểm & Main rotor & Tail rotor \\
\midrule
Độ lớn thrust/torque & Tính bằng BEMT 100 station & Tính bằng BEMT 100 station \\
Hướng lực & Thay đổi theo \(\vect{V}_b\), \(\vect{\omega}_b\), cyclic & Cố định bởi quay \(90^\circ\) quanh trục \(x\) \\
Đầu ra lực & \eqref{eq:main-force} & \([0,\ T_t,\ 0]^T\) \\
Đầu ra moment & \(\vect{r}_m \times \vect{F}_m + [0,0,-Q_m]^T\) & \(\vect{r}_t \times \vect{F}_t + [0,Q_t,0]^T\) \\
\bottomrule
\end{tabular}
\end{center}

\section{Kết luận}

\begin{enumerate}[leftmargin=1.8em]
  \item Cả main rotor và tail rotor đều lấy độ lớn thrust và torque từ cùng một
        lõi BEMT rời rạc 100 đoạn, thông qua \(C_T\) và \(C_Q\).
  \item Main rotor phức tạp hơn vì sau khi có độ lớn thrust, code còn tính hai
        góc nghiêng đĩa rotor từ \(\vect{V}_b\), \(\vect{\omega}_b\), cyclic
        pitch và các tham số rotor, rồi mới chiếu thrust vào hệ trục thân.
  \item Tail rotor đơn giản hơn: lực cục bộ luôn nằm dọc trục rotor, sau đó
        được quay cố định \(90^\circ\) quanh trục \(x\); vì vậy đầu ra lực của
        tail rotor luôn dọc trục \(+y\) của hệ thân trong mô hình hiện tại.
  \item Moment của cả hai rotor đều là tổng của moment đòn bẩy
        \(\vect{r}\times\vect{F}\) và moment cản khí động quanh trục rotor.
\end{enumerate}

\end{document}
